%% Career-Ops Cover Letter Template
%% Visual style matches templates/cv-template.tex (accent color, lmodern, ATS-friendly).
%% Build: pdflatex {file}.tex (or tectonic {file}.tex)

\documentclass[letterpaper,11pt]{article}

\usepackage[margin=0.85in]{geometry}
\usepackage[T1]{fontenc}
\usepackage[utf8]{inputenc}
\usepackage{lmodern}
\usepackage{microtype}
\usepackage{hyperref}
\usepackage{xcolor}
\usepackage{parskip}

\definecolor{accent}{HTML}{1F3A5F}

\hypersetup{
  colorlinks=true,
  urlcolor=accent,
  linkcolor=accent,
  pdftitle={ Ahnaf An Nafee - Cover Letter - Deepgram },
  pdfauthor={ Ahnaf An Nafee },
  pdfsubject={Cover Letter}
}

\input{glyphtounicode}
\pdfgentounicode=1

\pagestyle{empty}
\setlength{\parindent}{0pt}
\setlength{\parskip}{8pt}

\begin{document}

% --- Header (matches CV) ----------------------------------------------------
\begin{center}
  {\LARGE\bfseries\color{accent} Ahnaf An Nafee }\\[3pt]
  {\small Software Engineer \textbar{} AI Researcher \textbar{} DevOps \& Cloud Infrastructure }\\[4pt]
  \small
  \href{mailto:ahnafnafee@gmail.com}{ahnafnafee@gmail.com} \textbar{}
  Fairfax, VA\\
  \href{https://linkedin.com/in/ahnafnafee}{linkedin.com/in/ahnafnafee} \textbar{}
  \href{https://github.com/ahnafnafee}{github.com/ahnafnafee} \textbar{}
  \href{https://www.ahnafnafee.dev}{ahnafnafee.dev}
\end{center}

\vspace{20pt}

% --- Date -------------------------------------------------------------------
April 27, 2026

% --- Recipient block (multi-line; use \\ between lines) ---------------------
Andrew Seagraves\\Vice President of Research\\Deepgram

% --- Subject ----------------------------------------------------------------
\textbf{Re: Research Engineer, Machine Learning Systems}

% --- Greeting ---------------------------------------------------------------
Dear Andrew,

% --- Body (3 paragraphs separated by blank lines) ---------------------------
The cadence Deepgram is shipping research at --- Nova-3 STT going streaming-first, Aura-2 hitting sub-200ms TTFB --- only happens when training systems and research scientists move as one team, and that bridge is exactly where I want to live. I am a CS PhD candidate at George Mason's DCXR Lab working on generative AI for 3D content, with 2.5 years running production OpenShift, Kubernetes, and AWS at Paychex, plus a CTO tour at Dynasty 11 Studios shipping AI-adjacent consumer products. ``Research Engineer, Machine Learning Systems'' reads like the role written for someone who can engineer training pipelines on Monday and reason like a researcher on Tuesday, and that hybrid is the lane I have been building toward for years.

On the research side, my DCXR work centers on training generative models for 3D pipelines (UV mapping, neural rendering, NPR), which means I live the iterate-train-evaluate loop daily and know what research staff actually need from infrastructure. On the systems side at Paychex, I led cross-team standardization of Kubernetes/OpenShift deployment dictionaries across 10+ service teams, hardened the platform with RBAC, and automated the TLS lifecycle to eliminate certificate-related downtime by 100\%. Earlier, as CTO at Dynasty 11, I introduced GitHub Actions CI/CD that cut manual release effort by 85\% and migrated a monolith to microservices on AWS for an 80\% load reduction --- the kind of velocity a research org needs when experiments queue up faster than humans can ship them. I have not built ASR or TTS systems in production, but the training-infra primitives Deepgram uses --- job orchestration, experiment tracking, observability, data engineering for unstructured inputs --- are exactly the muscles I have spent my career strengthening, and the speech-domain ramp is something I am eager to close fast.

Deepgram's ``research that ships'' culture --- publications and patents on one side, a Nova/Aura production funnel on the other --- matches the trajectory I am on, and listening to your AIMinds episodes on context-conditioning for TTS made it feel less like a job description and more like a place built for the work I want to do. I would welcome a conversation about how your team scales the Nova-3 and Aura-2 training loops, and where a research engineer with this systems background could move the needle fastest.

% --- Closing ----------------------------------------------------------------
Thank you for your time and consideration.

\vspace{18pt}
Ahnaf An Nafee

\end{document}
